\section{有理数线}

\begin{proposition}{$\Q$是拓扑群}{}
	令$\mathfrak{F}=\left\{\left(-a,a\right)\subseteq\Q\middle|a\in\Q_{>0}\right\}$,那么$\mathfrak{F}$是$0$在$\Q$中的邻域基,它诱导的拓扑使$\Q$的加法群是拓扑群.
\end{proposition}

\begin{proposition}{$\Q$的第二可数性}{}
	$\left\{\left(-\frac{1}{n},\frac{1}{n}\right)\subseteq\Q\middle|n\geq 1\right\}$是$0$在有理数线中的可数邻域基,因此有理数线是第二可数空间.
\end{proposition}

\begin{corollary}{$\Q$中各点的邻域基的表示}{}
	对每个$x\in\Q$,集合$\left\{\left(x-a,x+a\right)\subseteq\Q\middle|a\in\Q_{>0}\right\}$和$\left\{\left(x-\frac{1}{n},x+\frac{1}{n}\right)\subseteq\Q\middle|n\geq 1\right\}$是$x$在$\Q$中的邻域基.
\end{corollary}

\begin{proposition}{$\Q$的Hausdorff性和非离散性}{}
	有理数线是非离散Hausdorff空间.
\end{proposition}

\begin{proposition}{$\Q$的邻里基}{}
	对每个$a\in\Q_{>0}$,令$U_{a}=\left\{\left(x,y\right)\in\Q\times\Q\middle|\left|x-y\right|<a\right\}$,那么$\left\{U_{a}\middle|a\in\Q_{>0}\right\}$和$\left\{U_{\frac{1}{n}}\middle|n\geq 1\right\}$都是有理数线的加法一致结构的邻里基.
\end{proposition}

\begin{proposition}{}{}
	\begin{enumerate}
		\item $\Q\to\Q_{>0},x\mapsto x^{+}$和$\Q\to\Q_{>0},x\mapsto x^{-}$都是一致连续的.
		\item $\Q\times\Q\to\Q,\left(x,y\right)\mapsto\sup\left(x,y\right)$和$\Q\times\Q\to\Q,\left(x,y\right)\mapsto\inf\left(x,y\right)$都是一致连续的.
	\end{enumerate}
\end{proposition}